\section[Tannakian 重构]{Tannakian 重构（3 月 11 日）}
\srcpage{38}

先说明一点：关于叠的那次课会被录下来，近期发布；想速成叠的入门、或复习本课用到的材料，
可以稍后异步观看。

今天讲 Tannakian 重构。口号是「交换性指示几何」：交换环是某个几何对象上的函数代数；
我们要把它范畴化。具体地，若 $R$ 是交换环，$\mathrm{Mod}_R$ 上有张量积
$M,N\mapsto M\otimes_R N$，结果仍是 $R$-模。这使 $\mathrm{Mod}_R$ 成为对称幺半范畴。

对称幺半范畴的定义有些技术性。我们关心的（与标准定义并不等价）是具备若干良好性质的阿贝尔范畴，
其中包括任意直和。Grothendieck 阿贝尔范畴满足这些运算。这样的范畴 $\cC$ 带对称幺半结构，
指的是给出行称张量积的函子 $\otimes:\cC\times\cC\to\cC$、称作单位的对象 $\mathbf{1}\in\cC$，
以及进一步的数据与条件。这些条件包括「与单位作张量是平凡的」，它由自然变换
\[
\mathbf{1}\otimes -\ \simeq\ -\otimes\mathbf{1}\ \simeq\ \id
\]
实现；还包括对应 $\otimes$ 的结合性、交换性（同构意义下）的自然变换，以及这些数据上的一些条件。
相应地有对称幺半函子的概念，它保持张量积与单位的数据。

要点是：$R$ 是交换环时，$\mathrm{Mod}_R$ 是以 $\otimes_R$ 为张量积、以 $R$ 为单位的对称幺半范畴。
以前我们只用 $R$-模做谱分解，现在把 $\mathrm{Mod}_R$ 看成范畴化的交换环。
这在 TFT 里有用：余维数升高，范畴的阶数也随之升高。

几何空间仍是 $X=\Spec R$，但我们要范畴化 $X$ 上的「函数环」，把 $R$ 换成
$\mathrm{Mod}_R=\QCoh(X)$。代数几何里大多数概形不是仿射的，这有点烦人，
也是与流形情形的重大差别。范畴化的一个好处是，更多的东西变成「仿射的」。
具体地，设 $X$ 是一般概形，有仿射开覆盖 $U=\{\Spec(R_i)\}$，则 $X$ 上拟凝聚层构成的范畴是
\[
\QCoh(X)=\varprojlim_{U}\ \mathrm{Mod}_{R_i},
\]
按你对 QC 的处理方式，这可以是定义也可以是定理。拟凝聚层以 $\otimes_{\cO_X}$ 为张量积、
以 $\cO_X$ 为单位，构成对称幺半范畴。此外，若 $f:X\to Y$ 是概形间的映射，
则有拉回函子 $f^*:\QCoh(Y)\to\QCoh(X)$，它是对称幺半的：$\cO_Y$ 拉回为 $\cO_X$，
张量积拉回为张量积。事实上这是万有的：$\QCoh(X)$ 满足如下万有性质——
只要有映射 $f:\Spec(R)\to X$，就有函子 $f^*:\QCoh(X)\to\mathrm{Mod}_R$。

总结一下：正如我们有 $\cO:\mathrm{Sch}^{\mathrm{op}}\to\mathrm{CommRing}$ 及良好的对应，
我们也有 $\QCoh:\mathrm{Sch}^{\mathrm{op}}\to\otimes\text{-}\mathrm{Cat}$。
此外还可推广到相当一般的叠，对叠 $X$ 定义 $\QCoh(X)$。允许哪些叠，稍后再说——
Artin 叠是好的例子，一般可以把它们想成概形对群作用的商。
\srcpage{39}

这给了我们对称幺半范畴的一大来源：还有别的获取途径，但这是我们的基地，
是最像几何、最亲近我们的那些范畴。于是可以问：给定任意对称幺半范畴，
能否把它重构为某个 $X$ 或 $R$ 上的 $\QCoh(X)$ 或 $\mathrm{Mod}_R$？

我们要构造一个右伴随，就像对 $\cO$ 所做的那样；这意味着范畴化的 $\Spec$ 满足
\[
\Hom_{\mathrm{Stack}}\bigl(X,\Spec\cC\bigr)
=\Hom_{\otimes\text{-}\mathrm{Cat}}\bigl(\cC,\QCoh(X)\bigr).
\]
归根结底，这是要求 $\Spec R\to\Spec\cC$ 的映射（前者是通常的环谱，后者是我们范畴化的概念）
与对称幺半函子 $\cC\to\mathrm{Mod}_R$ 自然地一一对应。\footnote{我们对态射的要求，
比「是对称幺半函子」略多一些。Tannakian 重构定理所需的精确条件，见 Lurie [Lur04]
或 Bhatt–Halpern-Leistner [BHL17]。}

$(\cO,\Spec)$ 伴随的单位（TODO：余单位？）给出任何概形的典范仿射化映射：
$X\mapsto\Spec\bigl(\cO(X)\bigr)$。例如 $X=\A^n$ 时，这就是通常的映射 $\A^n\to k$。
同样，上述伴随的单位（TODO：余单位？）给出「Tannakian 化」映射
$X\to\Spec\bigl(\QCoh(X)\bigr)$。

TODO：这里我漏讲了一点，但若 $S$ 是交换环、$\cC=\mathrm{Mod}_S$，则 $\cC$ 由单位生成，
意思是任何 $S$-模都有自由消解。也就是说，可以从 $S$ 出发，通过取直和、核与余核，
归纳地造出所有 $S$-模。也可以研究函子 $\Hom_{\mathrm{Mod}_S}(S,-):\mathrm{Mod}_S\to\mathrm{Ab}$，
它有一些很好的性质，同样刻画了「由单位生成」。这告诉我们（TODO：这一点我还没想通）
$\Spec(\mathrm{Mod}_S)=\Spec(S)$，这让人放心！

但这种范畴化的 $\Spec$ 真正的好处，是能理解许多非仿射的对象，甚至叠。
这些几何对象没有足够多的函数，却有足够多的拟凝聚层。概形上有大量层，例如摩天楼层，
尽管我们没有解范畴化后的类似物——$\delta$-函数。

\begin{theorem}[Tannakian 重构]
设 $X$ 是几何叠。则「Tannakian 化」映射 $X\to\Spec\QCoh(X)$ 是同构。
\end{theorem}

参考 Lurie [Lur04] 与 Bhatt–Halpern-Leistner [BHL17]。总之，几何叠是「范畴意义下仿射的」；
这个定理的力量在于，你合理能遇到的东西几乎都是几何叠。例如所有拟紧概形都是。
几何叠的完整定义是：具有仿射对角的拟紧叠。这在实践中相当具体：若 $X$ 是这样一个叠，
可以取仿射簇 $U$ 与作用其上的仿射群胚 $G$ 来构造它，此时 $X\cong U/\!/G$。

\begin{remark}
「范畴化的 $\delta$-函数」这件事并不太令人惊讶：层似乎同时范畴化了函数与分布/测度。
这与层既能拉回又能前推有关。
\end{remark}

接下来的时间用来讲例子。

\begin{example}
作为第一个例子，设 $X$ 是拟紧概形。我们可以通过拉回函子 $f^*:\QCoh(X)\to\mathrm{Mod}_R$
来探测映射 $\Spec R\to X$，前面的 Tannakian 重构定理告诉我们这两者是等价的。例如 $X$ 的 $k$-点即映射
$\Spec k\to X$，等价于拉回 $f^*:\QCoh(X)\to\QCoh(\Spec k)=\Vect_k$。
也就是说，一个点是到 $\Vect$ 的张量函子。
\end{example}

典型的例子在表示论里。设 $G$ 是群，则 $\Rep G$（在选定的域 $k$ 上考虑表示）
是对称幺半范畴：单位是平凡的一维表示，张量积是向量空间的张量积带上对角的 $G$-作用。
这是自然地依赖于 $G$ 的对称幺半范畴。

现在把它塞进 Tannakian 的机器：$\Spec(\Rep G)$ 是某个东西上拟凝聚层的范畴吗？
要想这件事，需要多一点数据。表示是向量空间加上 $G$-作用，所以可以忘掉作用，
得到忠实对称幺半函子 $\Rep(G)\to\Vect$。一般地，到 $\Vect$ 的忠实对称幺半函子叫
纤维函子（fiber functor）\yizhu{「纤维函子」之名来自 Tannaka 对偶：当 $\cC=\Rep(G)$ 时，它就是把表示送到底层向量空间、即「在一点处取纤维」的函子。}{B（nLab "Tannaka duality"）}。它是特殊的：对簇 $X$，$\QCoh(X)$ 永远没有这样的函子，
因为这等价于一个映射 $i:\Spec k\to X$，而沿它拉回永远不是忠实的，除非 $X$ 本身就是点。
\srcpage{40}

所以有纤维函子与「是概形」多少正交。可以把它理解为：$\Spec\cC$ 允许一个从点出发的
忠实平坦映射，直观地说就是 $\Spec\cC$ 是点的商。这里「叠性」很重要，
唯其如此才可能非平凡。

特别地（这正是经典意义上的 Tannakian 重构），若 $G$ 是仿射代数群，
即仿射簇范畴中的群对象（于是 $G$ 是仿射簇，乘法与求逆都是簇的映射，如 $\mathrm{GL}_n(k)$），
则 $\Spec(\Rep G)=BG=\mathrm{pt}/G$，而纤维函子等同于商映射 $\mathrm{pt}\to\mathrm{pt}/G$。
我们可以把 $G$ 恢复成我们所施加的那些关系。

通常的说法是：给定 $\cC$ 与纤维函子 $f$，可以在 $\Aut(f)$ 上构造仿射代数群的结构，
并且 $\cC\simeq\Rep G$。

世上有很多怪群，没有有限维代数表示，即没有到 $\mathrm{GL}_n(k)$ 的映射。
从这个视角看不到它们；一般地 $\Spec(\cC)$ 是 $G$ 的投射代数完备化（proalgebraic
completion），是为使其更「代数」而必须做的一些技术处理。某种意义上，
投射代数完备化是由 $G$ 能造出的万有代数群。

下面到自然界里寻找带纤维函子的对称幺半范畴。（例子会偏向代数几何，
因为我们的测试范畴是 $R$-模范畴；若你操作的是 $C^*$-代数上的模之类，例子会更有拓扑味。）
我们可以通过造带纤维函子的范畴——即带结构的向量空间——来造群。

\begin{example}
考虑 $\Z$-分次向量空间构成的范畴；忘掉分次给出到 $\Vect$ 的忠实忘却函子。
分次张量积为
\[
(M\otimes N)_n=\bigoplus_{i+j=n}M_i\otimes N_j .
\]
这里不用 Koszul 符号规则。因此忘却函子是对称幺半的，Tannakian 重构定理告诉我们，
存在群 $G$ 使 $\Vect^{\Z}\simeq\Rep G$。此例中我们知道答案：$\Z$-分次向量空间
等同于乘法群 $\G_m$ 的表示（取等型分量）；不可约 $\G_m$-表示由 $n\in\Z$ 标号，
$z\in\G_m$ 在 $k$ 上作用为 $z^n$。这又是谱分解/Pontrjagin 对偶的一种形式：
$\Z$-分次与 $\G_m$-作用对偶。所以 $\Vect^{\Z}\simeq\QCoh(\G_m)$。
\end{example}

Pontrjagin 对偶说：$G$ 是紧群时其对偶 $G^\vee$ 离散，且紧群的表示范畴是半单的；
在代数情形，表示范畴半单意味着群约化。$G$ 约化当且仅当 $\QCoh(\mathrm{pt}/G)=\Rep G$
半单。Pontrjagin 对偶的相当一部分故事，可以从紧群搬到这个场景。

\begin{example}
这个例子是「纤维函子」这个名字的由来。设 $X$ 是连通（且局部道路连通）的拓扑空间，
考虑 $X$ 上向量空间局部系统构成的对称幺半范畴 $\mathrm{Loc}(X)$；
对称幺半结构是逐点张量积。取基点 $x\in X$，在 $x$ 处取纤维给出函子 $\mathrm{Loc}(X)\to\Vect$；
因为张量积是逐点的，这是对称幺半函子；因为 $X$ 连通，它忠实。所以取纤维就是纤维函子。
于是得到群 $G$ 使 $\mathrm{Loc}(X)=\Rep G$。这个 $G$ 本质上是 $\pi_1(X,x)$，
因为 $\pi_1$ 在 $k$-向量空间上的表示，通过单值表示等同于 $X$ 上 $k$-向量空间的局部系统。
然而任何群都可能是 $\pi_1$，但不是所有群都有有限维代数表示——这就是上面「本质上」的意思：
一般地，得到的 $G$ 是 $\pi_1(X)$ 的投射代数完备化。
\end{example}

\begin{example}
设 $k$ 是域，考虑 $\cC:=\Vect_k$。取 $k$ 的代数闭包 $\bar{k}$；则 $-\otimes_k\bar{k}$
是忠实对称幺半函子，把 $\cC$ 等同于某个群 $G$ 的 $\Rep G$，事实上 $G=\Gal(\bar{k}/k)$。
这是 Galois 下降的一种形式：Galois 群来自域扩张的自同构群，我们在这里看到了它。
\end{example}

尽管我们立足于代数几何，这些例子却来自别处，而我们给它们配上了代数几何的对象。
这个想法我们以后还会见到。

下次把它与 Fourier–Mukai 理论结合起来：交换环的一个好来源是阿贝尔群 $G$ 的群代数。
我们将有由乘法给出的前推
\[
\mu_*:\mathrm{Sh}(G)\times\mathrm{Sh}(G)\longrightarrow\mathrm{Sh}(G),
\]
即卷积。这对任何群都有意义，并在 $G$ 阿贝尔时给出对称幺半结构。
我们没说明这里出现的是哪一种层——既然目的是得到对称幺半范畴，拟凝聚与否无所谓。
\srcpage{41}

于是可以定义 $G$ 的 Fourier–Mukai 对偶或 1-移位 Cartier 对偶为
\[
G^\vee[1]:=\Spec\bigl(\mathrm{Sh}(G),*\bigr).
\]
例如 $G=\Z$ 时，$\Z$ 上的层就是 $\Z$-分次向量空间，卷积就是分次张量积。
所以这是范畴化的群代数，其 1-移位 Cartier 对偶是 $B\G_m$。这里移位是显然的：
$\Z$ 的通常 Cartier 对偶是 $\G_m$，但我们移位之后得到的却是 $B\G_m$，
尽管我们走的是一条类似的推理路线，只是把它范畴化了。

下次我们更多地讨论 1-移位 Cartier 对偶，并让它看起来很像 Fourier 理论。
